\color{unidarkgreen}{\subsection{Газовая защита}}
\color{black}

\subsubsection{Общие сведения}

\begin{enumerate}[label=\arabic{section}.\arabic{subsection}.\arabic{subsubsection}.\arabic*,  align=left, labelsep=4pt, leftmargin=0pt, itemindent=57pt]

\item
Газовая защита трансформатора предназначена для защиты от повреждений внутри кожуха,
сопровождающихся выделением газа, и понижением уровня масла. 

\begin{figure}[H]
\centering
\includegraphics[width=0.82\textwidth,height=0.82\textheight,keepaspectratio]{\fbpath/011.1071 ГЗ Т/img/GZT_sign.pdf}
\captionof{figure}{Функциональный блок ГЗ Т сигнальная ступень}\label{pic:GZT_sign}
\end{figure}

\begin{figure}[H]
\centering
\includegraphics[width=0.82\textwidth,height=0.82\textheight,keepaspectratio]{\fbpath/011.1071 ГЗ Т/img/GZT_otkl.pdf}
\captionof{figure}{Функциональный блок ГЗ Т отключающая ступень}\label{pic:GZT_otkl}
\end{figure}

\item
Предусмотрена возможность приема сигналов срабатывания сигнального контакта ГЗ Т,
отключающего контакта ГЗ Т и низкой изоляции цепи отключающего контакта ГЗ Т, по разным цепям от
двух источников ПДС1 и ПДС2. Для приема сигналов от ПДС1 предусмотрены входные сигналы \textit{«Ср.
сигн контакта ГЗ Т 1 цепь»}, \textit{«Ср. откл контакта ГЗ Т 1 цепь»} и \textit{«Низкая изол. ГЗ Т откл 1 цепь»}, а для
приема сигналов от ПДС2 предусмотрены входные сигналы \textit{«Ср. сигн контакта ГЗ Т 2 цепь»}, \textit{«Ср. откл
контакта ГЗ Т 2 цепь»} и \textit{«Низкая изол. ГЗ Т откл 2 цепь»}. Если сигналы срабатывания ГЗ Т заводятся
непосредственно со шкафа ШОМ силового трансформатора конфигурируются только входные сигналы
\textit{«Ср. сигн контакта ГЗ Т 1 цепь»}, \textit{«Ср. откл контакта ГЗ Т 1 цепь»} и \textit{«Низкая изол. ГЗ Т откл 1 цепь»}, а
действие ГЗ Т по второй цепи сигнального и отключающего контакта выводиться из работы с помощью
уставок \textbf{«ГЗ Т сигн 1 цепь»} и \textbf{«ГЗ Т откл 2 цепь»}.

\item
Ввод в работу ГЗ Т осуществляется с помощью уставок \textbf{«ГЗ Т сигн 1 цепь»}, \textbf{«ГЗ Т сигн 2 цепь»},
\textbf{«ГЗ Т откл 1 цепь»} и \textbf{«ГЗ Т откл 2 цепь»}. Оперативный перевод ГЗ Т на сигнал производится с помощью
команд управления \textit{«Перевод ГЗ Т откл на сигнал»}. Предусматривается возможность действия
срабатывания сигнального контакта ГЗ Т  на отключение трансформатора с помощью уставки \textbf{«Откл Т
от ГЗ сигн»}.

\item
Оперативный ввод/вывод ГЗ Т сигнальной ступени производится с помощью БУ (блока управления
функциями) \textit{«Вывод ГЗ Т сигн»}, а оперативный ввод/вывод ГЗ Т отключающий ступени производится с
помощью БУ (блока управления функциями) \textit{«Вывод ГЗ Т откл»}.

\item
Непрерывный контроль изоляции цепи отключающего контакта ГЗ Т осуществляется с помощью
реле контроля изоляции (РКИ). В случае снижения сопротивления изоляции РКИ подает сигнал \textit{«Низкая
изол. ГЗ Т откл 1 цепь»} (\textit{«Низкая изол. ГЗ Т откл 2 цепь»}) для блокировки срабатывания ГЗ Т от
отключающего контакта. Время срабатывания блокировки от сигналов РКИ задается с помощью уставки
\textbf{«Тбл ГЗ Т откл»}. Ввод блокировки ГЗ Т от РКИ осуществляется с помощью уставки \textbf{«Блок ГЗ Т откл»}.
Для сброса блокировки от РКИ, используется команда управления \textit{«Возврат ГЗ Т откл»}.

\item
ГЗ Т аналогична ГЗ РПН, ГЗ ЛРТ.

\small
\begin{longtable}{|>{\centering\arraybackslash}m{5.3cm}|>{\centering\arraybackslash}m{4.2cm}|>{\centering\arraybackslash}m{1cm}|>{\centering\arraybackslash}m{3.3cm}|>{\centering\arraybackslash}m{1.8cm}|}
% Первая страница: полный заголовок
\caption{Перечень уставочных параметров ГЗ\hfill\vspace{-0.5\baselineskip}}\label{desc:tbl1}\\
\hline
\rowcolor{gray!30}
Наименование уставки & Диапазон & Шаг & Значение по умолчанию & Единица измерения \\
\hline
\endfirsthead

% Промежуточные страницы: «Продолжение таблицы»
\caption*{\hspace{3pt}\emph{Продолжение таблицы \ref{desc:tbl1}}\hfill\ }\vspace{-0.5\baselineskip} \\
\hline
\rowcolor{gray!30}
Наименование уставки & Диапазон & Шаг & Значение по умолчанию & Единица измерения \\
\hline
\endhead

%===m>ДТЗ|-
%===m
\end{longtable}
\normalsize	

\end{enumerate}